%==============================================================================
% CHAPITRE 6 — Déploiement (en ligne uniquement)
% URL application : https://bts-dashboard-tt.streamlit.app/
% Prérequis : \usepackage{tikz}, \usetikzlibrary{positioning,arrows.meta}
%             \usepackage{booktabs}, \usepackage{hyperref}
%==============================================================================

\chapter{Déploiement et mise en exploitation du tableau de bord BTS~EMS}
\label{chap:deploiement}

%------------------------------------------------------------------------------
\section{Introduction}
\label{sec:dep-intro}
%------------------------------------------------------------------------------

Les chapitres précédents ont produit trois livrables analytiques~: un modèle
LightGBM de prédiction énergétique~(NB1), un ensemble de détecteurs
d'anomalies~(NB2) et un moteur de décision par apprentissage par renforcement~(NB3).
Ces sorties ne génèrent de valeur opérationnelle que mises à disposition d'un
exploitant réseau dans un outil accessible et filtrable.

Ce chapitre documente le \textbf{déploiement en ligne} du tableau de bord
\textbf{BTS~EMS}, accessible à l'adresse suivante~:
\begin{center}
  \url{https://bts-dashboard-tt.streamlit.app/}
\end{center}

Le périmètre retenu pour ce projet est exclusivement ce déploiement
\textbf{cloud}~: application hébergée sur Streamlit Community Cloud, artefacts
machine learning servis depuis le Hub Hugging~Face \texttt{molkab/dashboard},
persistance des comptes et des alertes en SQLite sur l'infrastructure Streamlit.
Aucune installation locale des notebooks, des dossiers de sortie (\texttt{VF/})
ou de l'application n'est requise pour l'utilisateur final. La conteneurisation
(Docker, reverse proxy dédié) reste une perspective d'industrialisation chez
Tunisie Telecom, hors livrable de ce stage.

% CAPTURE : page d'accueil ou barre d'URL Streamlit
% \includegraphics[width=0.9\textwidth]{figures/fig6_0_url_streamlit.pdf}

%------------------------------------------------------------------------------
\section{Architecture globale de la solution déployée}
\label{sec:dep-archi}
%------------------------------------------------------------------------------

\subsection{Vue d'ensemble}
\label{subsec:dep-vue-ensemble}

La solution déployée s'organise en quatre couches, représentées sur la
figure~\ref{fig:archi-deploiement}.

\begin{enumerate}
  \item \textbf{Couche analytique (hors ligne, amont).}
    Les notebooks \texttt{nb1-prediction.ipynb}, \texttt{nb2-d.ipynb} et
    \texttt{nb3-p.ipynb} produisent modèles (\texttt{.joblib}), métriques
    (\texttt{.json}) et tables (\texttt{.parquet}). Cette étape est réalisée
    en amont du déploiement~; les fichiers sont ensuite publiés sur le Hub.

  \item \textbf{Couche de publication en ligne.}
    L'ensemble des artefacts est centralisé sur le dépôt Hugging~Face
    \texttt{molkab/dashboard} (\texttt{HF\_REPO\_ID}). Au démarrage de
    l'application, \texttt{artifact\_path()} télécharge ou met en cache les
    fichiers nécessaires (\texttt{USE\_HF\_HUB=True}).

  \item \textbf{Couche applicative (serveur Streamlit).}
    \texttt{services/data\_service.py} charge les artefacts depuis le Hub,
    fusionne les jeux NB2 et NB3 (\texttt{enrich\_dashboard\_data}) et calcule
    les indicateurs (\texttt{compute\_filtered\_kpis},
    \texttt{harmonize\_nb3\_economies}).

  \item \textbf{Couche de présentation.}
    Sept entrées de menu pour l'administrateur (six pour l'ingénieur), accessibles
    via navigateur~; comptes et alertes en SQLite côté hébergeur.
\end{enumerate}

Le tableau de bord \emph{relit} et \emph{agrège} des résultats pré-calculés par
les notebooks. Il ne ré-exécute pas l'inférence LightGBM à chaque filtre de la
barre latérale. Seule la page \textbf{Simulation} relance, à chaque tick simulé,
une chaîne proche du pipeline NB3 (\texttt{run\_nb\_pipeline}, repli
\texttt{apply\_offline\_nb23} si un modèle Hub est momentanément indisponible).

\begin{figure}[htbp]
  \centering
  \begin{tikzpicture}[
      node distance = 0.65cm and 1.1cm,
      box/.style   = {rectangle, draw, rounded corners=3pt,
                      minimum width=3.5cm, minimum height=0.65cm,
                      font=\small, align=center},
      arrow/.style = {->, >=stealth, thick},
      cloud/.style = {box, fill=blue!8},
      store/.style = {box, fill=orange!12},
      app/.style   = {box, fill=green!10},
      ui/.style    = {box, fill=purple!10},
    ]
    \node[cloud] (nb) {Notebooks NB1--NB3\\(amont)};
    \node[store, right=of nb] (hub)
      {Hub HF\\\texttt{molkab/dashboard}};
    \node[app, right=of hub] (cloud)
      {Streamlit Cloud\\cache + SQLite};
    \node[app, right=of cloud] (svc)
      {\texttt{data\_service.py}\\enrich + KPI};
    \node[ui, right=of svc] (web)
      {\url{bts-dashboard-tt}\\.streamlit.app}};
    \node[ui, right=of web] (users) {Admin /\\Ingénieur};
    \draw[arrow] (nb) -- (hub);
    \draw[arrow] (hub) -- (cloud);
    \draw[arrow] (cloud) -- (svc);
    \draw[arrow] (svc) -- (web);
    \draw[arrow] (web) -- (users);
  \end{tikzpicture}
  \caption{Architecture de déploiement en ligne~: artefacts sur Hugging~Face,
           application sur Streamlit Cloud, accès navigateur.}
  \label{fig:archi-deploiement}
\end{figure}

%------------------------------------------------------------------------------
\section{Gestion des artefacts et intégration des notebooks}
\label{sec:dep-artefacts}
%------------------------------------------------------------------------------

\subsection{Registre et chargement depuis le Hub}
\label{subsec:dep-registre}

Le module \texttt{services/data\_service.py} référence chaque fichier attendu
dans \texttt{ARTIFACT\_REGISTRY} (notebook d'origine, type, usage interface).
La fonction \texttt{artifact\_path(filename)} résout l'URL Hub et le cache
local du serveur Streamlit (\texttt{.cache/huggingface}), sans dépendre de
copies manuelles sur le poste de l'utilisateur. Le tableau~\ref{tab:artefacts}
résume les artefacts principaux.

\begin{table}[htbp]
  \centering
  \caption{Artefacts Hub utilisés par le tableau de bord en ligne}
  \label{tab:artefacts}
  \begin{tabular}{@{}p{0.7cm}p{5.5cm}p{6cm}@{}}
    \toprule
    \textbf{NB} & \textbf{Fichiers clés (Hub)} & \textbf{Usage interface} \\
    \midrule
    NB1 &
    \texttt{df\_full\_processed.parquet},
    \texttt{resultats\_modeles.json},
    \texttt{quantile\_models.joblib} &
    Courbes réel vs prédit, bandes Q10--Q90, R² \\
    \addlinespace
    NB2 &
    \texttt{df\_avec\_anomalies.parquet},
    \texttt{resultats\_anomalie.json} &
    Scores, votes, seuil d'ensemble (ou dérivé) \\
    \addlinespace
    NB3 &
    \texttt{streamlit\_data.parquet},
    \texttt{decisions\_par\_station.parquet},
    \texttt{rapport\_optimisation.json} &
    KPI, actions, économies, simulation \\
    \bottomrule
  \end{tabular}
\end{table}

\subsection{Enrichissement et harmonisation}
\label{subsec:dep-enrich}

\texttt{enrich\_dashboard\_data()} fusionne le jeu NB3 avec les colonnes NB2 sur
\texttt{(station\_id, timestamp)}. \texttt{harmonize\_nb3\_economies()} applique~:
\begin{equation}
  \text{economie\_kwh}
  = \min\!\bigl(\max(\text{règles},\; \text{RL}),\;
               0{,}48 \times \text{consommation\_kwh}\bigr)
  \label{eq:eco}
\end{equation}

Si \texttt{seuil\_ensemble} est absent de \texttt{resultats\_anomalie.json}, le
seuil est dérivé du quantile des scores dans \texttt{df\_avec\_anomalies.parquet},
calibré sur \texttt{pct\_anomalies} de \texttt{kpi\_reseau.json}.

%------------------------------------------------------------------------------
\section{Structure de l'application Streamlit en ligne}
\label{sec:dep-streamlit}
%------------------------------------------------------------------------------

\subsection{Hébergement et accès}
\label{subsec:dep-hebergement}

L'application est déployée sur \textbf{Streamlit Community Cloud}. L'URL
publique \url{https://bts-dashboard-tt.streamlit.app/} pointe vers le point
d'entrée \texttt{app.py} du projet~: configuration de page, chargement des
secrets plateforme (\texttt{SECRET\_KEY}, \texttt{HF\_TOKEN}, SMTP), puis
routeur \texttt{ui/dashboard.py}. L'utilisateur n'exécute aucune commande
locale~; il se connecte via navigateur web (authentification requise).

% CAPTURE : page connexion Streamlit Cloud
% figures/fig6_7_login.pdf

\subsection{Écrans et correspondance avec les notebooks}
\label{subsec:dep-ecrans}

\begin{table}[htbp]
  \centering
  \caption{Écrans du tableau de bord en ligne et notebooks associés}
  \label{tab:ecrans}
  \begin{tabular}{@{}p{2.6cm}p{1cm}p{9.2cm}@{}}
    \toprule
    \textbf{Écran} & \textbf{NB} & \textbf{Contenu principal} \\
    \midrule
    Accueil        & --- & KPI, camembert modes, top~5 critiques \\
    Carte          & --- & Carte Folium, actions NB3 \\
    Prédiction     & NB1 & Réel / prédit / Q10--Q90, SHAP (admin) \\
    Anomalies      & NB2 & Scores, alertes, détecteurs (admin) \\
    Optimisation   & NB3 & Économies retenues, actions, gains expert/RL \\
    Simulation     & NB3 & Replay horaire, journal d'alertes \\
    Configuration  & --- & Stations, import dataset, utilisateurs (admin, 3 onglets) \\
    \bottomrule
  \end{tabular}
\end{table}

% CAPTURE : menu latéral 7 entrées — fig6_3_menu_admin.pdf
% CAPTURE : Accueil KPI + camembert — fig6_4_accueil.pdf
% CAPTURE : Optimisation — fig6_5_optimisation.pdf

Les filtres globaux alimentent \texttt{compute\_filtered\_kpis()} sans
réentraînement des modèles.

%------------------------------------------------------------------------------
\section{Sécurité, configuration et persistance}
\label{sec:dep-securite}
%------------------------------------------------------------------------------

\subsection{Authentification et rôles}
\label{subsec:dep-auth}

Authentification via \texttt{auth\_service.py} et \texttt{utils/security.py}
(bcrypt, SQLite sur le serveur Streamlit). Verrouillage après cinq échecs de
connexion (cinq minutes).

\begin{itemize}
  \item \textbf{Administrateur}~: sept écrans, KPI en DT, détails ML.
  \item \textbf{Ingénieur}~: six écrans, stations assignées via
    \texttt{get\_user\_stations()} dans \texttt{data\_service.py} (RLS).
\end{itemize}

\subsection{Protections applicatives}
\label{subsec:dep-middleware}

\texttt{security/middleware.enforce()} est appelé en tête de chaque page~:
rate limiting ($\approx$200 événements/min), expiration de session (1~h), CSRF
en session. Streamlit ne pose pas d'en-têtes HTTP (CSP, HSTS) comme un reverse
proxy Nginx~; en production institutionnelle, un proxy devrait compléter ces
aspects devant l'URL publique.

\subsection{Configuration cloud}
\label{subsec:dep-config}

Constantes métier dans \texttt{config/settings.py}~: \texttt{PRIX\_KWH\_TN},
\texttt{QOS\_SEUIL\_DEFAULT}, \texttt{NB3\_MAX\_ECO\_FRAC}. Les secrets sont
définis dans l'espace \emph{Secrets} de Streamlit Cloud (équivalent de
\texttt{.env}, non exposé aux utilisateurs).

%------------------------------------------------------------------------------
\section{Déploiement retenu}
\label{sec:dep-modes}
%------------------------------------------------------------------------------

Le tableau~\ref{tab:modes-deploy} résume le mode de déploiement effectif.

\begin{table}[htbp]
  \centering
  \caption{Déploiement de l'application BTS~EMS}
  \label{tab:modes-deploy}
  \begin{tabular}{@{}p{3.2cm}p{10.8cm}@{}}
    \toprule
    \textbf{Composant} & \textbf{Description} \\
    \midrule
    \textbf{Interface utilisateur}
      & \url{https://bts-dashboard-tt.streamlit.app/} — Streamlit Community Cloud,
        Python~3.11 (\texttt{runtime.txt}) \\
    \addlinespace
    \textbf{Artefacts ML}
      & Hub Hugging Face \texttt{molkab/dashboard} — téléchargement et cache
        serveur à chaque session \\
    \addlinespace
    \textbf{Données applicatives}
      & SQLite (comptes, alertes, paramètres) sur l'hébergeur Streamlit \\
    \addlinespace
    \textbf{Non livré}
      & Installation locale, dossiers \texttt{VF/}, Docker, Nginx, sonde
        \texttt{health.py} dédiée \\
    \bottomrule
  \end{tabular}
\end{table}

\paragraph{Perspective.}
Une industrialisation chez Tunisie Telecom pourrait migrer l'URL vers un domaine
interne, ajouter un reverse proxy TLS et planifier la republication des
artefacts sur le Hub après chaque cycle d'entraînement des notebooks.

%------------------------------------------------------------------------------
\section{Page Simulation et inférence}
\label{sec:dep-simulation}
%------------------------------------------------------------------------------

La page \textbf{Simulation} charge des profils via \texttt{synthetic\_bts.py}
(calendrier \texttt{calendar\_tn.py}) et appelle \texttt{sim\_inference.py}
(pipeline Hub ou repli léger). Il s'agit d'un scénario \emph{what-if}, pas d'une
télémétrie live du réseau.

% CAPTURE : Simulation — fig6_6_simulation.pdf

%------------------------------------------------------------------------------
\section{Validation logicielle et limites}
\label{sec:dep-tests}
%------------------------------------------------------------------------------

\subsection{Tests automatisés}
\label{subsec:dep-tests-suite}

Une suite \texttt{pytest} (19 tests) valide la logique métier hors navigateur
(tableau~\ref{tab:tests}). Elle est exécutée lors de l'intégration continue du
code source~; elle ne remplace pas un test utilisateur sur l'URL de production.

\begin{table}[htbp]
  \centering
  \caption{Tests automatisés (intégration continue)}
  \label{tab:tests}
  \begin{tabular}{@{}p{7cm}p{7cm}@{}}
    \toprule
    \textbf{Fichier} & \textbf{Objet} \\
    \midrule
    \texttt{test\_nb\_metrics.py} & Économies, eq.~\eqref{eq:eco} \\
    \texttt{test\_decision\_service.py} & Modes NB3 \\
    \texttt{test\_security\_password.py} & Bcrypt \\
    \texttt{test\_middleware\_rate\_limit.py} & Rate limit login \\
    \texttt{test\_dashboard\_smoke.py} & Menu 7/6 entrées \\
    \bottomrule
  \end{tabular}
\end{table}

% CAPTURE : pytest 19 passed — fig6_8_pytest.pdf

\subsection{Limites}
\label{subsec:dep-limites}

\begin{enumerate}
  \item Dépendance à la disponibilité du Hub Hugging Face et du cache serveur.
  \item Métriques R² NB1 figées (JSON global, non recalculées sur filtre).
  \item Seuil anomalie parfois dérivé du parquet.
  \item Simulation = aide à la décision, pas jumeau temps réel du réseau.
\end{enumerate}

%------------------------------------------------------------------------------
\section{Conclusion}
\label{sec:dep-conclusion}
%------------------------------------------------------------------------------

Ce chapitre a présenté le déploiement \textbf{en ligne} du tableau de bord
BTS~EMS sur \url{https://bts-dashboard-tt.streamlit.app/}, alimenté par les
artefacts NB1--NB3 du Hub \texttt{molkab/dashboard}. L'architecture en quatre
couches (notebooks amont, Hub, serveur Streamlit, navigateur) assure une
traçabilité claire sans installation locale pour l'exploitant. Les mécanismes
de sécurité (bcrypt, rate limiting, RLS) et les 19 tests pytest complètent ce
déploiement cloud, ouvert à une industrialisation ultérieure chez Tunisie Telecom.
